\subsection{Omnidirectional Base}\label{sec:omnibase}
The Dashgo B1 used until 2026 had a differential drive, so aligning with a table or shelf took several turn-and-drive corrections before every pick or place. The team built its own holonomic base instead; its firmware\footnote{\url{https://github.com/RoBorregos/home-custom-base}} is open source.

\subsubsection{Mechanical Design.}
A frame of 4040 aluminum profiles carries a central column that raises the arm. Mecanum wheels keep their axles parallel to the frame, so the motors bolt directly to the profiles without the angled mounts that omni wheels need~\cite{taheri2015kinematic}. Each is driven through a 3D-printed 9:1 gearbox, a low ratio that lets the robot be pushed by hand when idle. The drives cap wheel speed at 1.33\,m/s and acceleration at about 1.3\,m/s$^2$, and Nav2 limits the base to 0.7\,m/s along both axes and 1.5\,rad/s in yaw. Printing let the team iterate gear materials cheaply: PLA and rigid-resin gears broke under load, and the current ABS gears hold but wear. The next step is carbon-fiber nylon (PAHT-CF), whose heat deflection temperature (about 190\,$^\circ$C against 87\,$^\circ$C for ABS) and fiber stiffness should reduce tooth wear and deformation.

\subsubsection{Electronics.}
Each wheel has an ODrive S1 motor driver whose absolute encoder lets it start in closed loop at power-on without a calibration motion; braking energy goes to a resistor instead of the battery that also feeds the computers. A team-designed shield connects an STM32H755 microcontroller to the 500\,kbps CAN bus, the IMU and the gamepad radio. Because drive feedback can delay setpoints on the shared bus, automatic retransmission stays on, and the CAN bit timing tolerates the largest clock mismatch, since a stricter setting repeatedly forced a drive off the bus.

\subsubsection{Firmware and Odometry.}
FreeRTOS tasks separate commands, drive control with mecanum inverse kinematics~\cite{taheri2015kinematic}, IMU and telemetry. If no velocity command arrives for 300\,ms, the firmware commands zero velocity, and each drive's own 0.5\,s watchdog stops the wheel even if the microcontroller or the bus fails. CAN and serial errors are recovered without a reboot, and a drive that drops out of closed-loop control is re-enabled at zero velocity. The normally closed e-stop also trips on a broken wire, and its release needs a hold time so a bouncing contact cannot re-engage the motors. Tasks queue their errors, which a ROS~2 dashboard displays. On the microcontroller, a six-state extended Kalman filter (EKF) fuses wheel velocities with IMU yaw and yaw rate at 50\,Hz, so odometry exists from boot and survives a ROS restart. Since mecanum rollers slip, the filter relies on the IMU for rotation, falling back to the wheels if it stops.